\chapter{Simulating Infectious Diseases}

\section{Euler's method}

Recall $f(x+\Delta x) \approx f(x)+f'(x)\Delta x$. \emph{Euler's
method} uses this repeatedly for $\dot x = f(x,t)$:  
\[
x_{t+\Delta t} \approx x_t + f(x_t,t)\,\Delta t,
\]
then repeat. For an SIR system, $x_t = (S,I,R)$:
\begin{lstlisting}
import numpy as np

def euler(beta, gamma, dt, n):
    S, I = np.zeros(n), np.zeros(n)
    S[0], I[0] = 0.99, 0.01
    for t in range(n - 1):
        dS =  - beta * S[t] * I[t]
        dI = beta * S[t] * I[t] - gamma * I[t]

        # Euler's method: new state = old state + rate * step size
        S[t + 1] = S[t] + dS * dt
        I[t + 1] = I[t] + dI * dt
        
    return S, I
\end{lstlisting}

\section{Incidence vs.~prevalence}
$I_t$ is the per-capita \emph{prevalence}: the proportion infected at
a time $t$. Surveillance rarely measures prevalence. Instead,
surveillance efforts reoirt \emph{incidence}: the number of \emph{new}
infections over a reporting interval, e.g.\ a week:

\begin{definition}
The \emph{incidence} over an interval $[t-\Delta t,\,t]$ is the number
of new infections:
\[
\text{Incidence}(t) = N \int_{t-\Delta t}^{t} \lambda_s\, S_s \, ds.
\]
\end{definition}
Tyically we normalize to find the \emph{per-capita} incidence

\begin{lstlisting}
import numpy as np

def incidence_euler(beta, gamma, dt, k, n):
    S, I = 0.99, 0.01
    incidence = np.zeros(n)
 
    for t in range(n-1):
        for s in range(k):
            new_infections = beta * S * I
            dS =  - new_infections
            dI = new_infections - gamma * I

            # Euler's method: new state = old state + rate * step size
            S += dS * dt
            I +=  dI * dt

            # Accumulate new infections
            incidence[t+1] += new_infections * dt 
        
    return incidence
\end{lstlisting}

Note: since the Euler time step (e.g.~day) may differ from the
sampling rate\footnote{To be precise, the Euler step may differ from
the \emph{inverse} of the sampling rate.} (e.g.~one week) we had to
implement another loop that rubs over $0,...,k$.

\section{Log-transformations}

Working with $\log$ variables is more stable numeically. We can easily
find that

\begin{align*}
  \dot{\log(S)} &= \frac{1}{S} \dot{S}  \\
  &=  \frac{-\beta S I}{S} \\
  %&= -\beta I \\
  &= -\beta e^{\log(I)}\\ 
  \dot {\log(I)} &= \frac{\beta SI -\gamma I}{I} \\
  &= -\beta e^{\log(S)} - \gamma.
\end{align*}

\begin{lstlisting}
import numpy as np

def log_euler(beta, gamma, dt, n):
    logS, logI = np.zeros(n), np.zeros(n)
    logS[0], logI[0] = np.log(0.99), np.log(0.01) 
    for t in range(n - 1):
        dlogS =  - beta * np.exp(logI[t])
        dlogI = beta * np.exp(logS[t]) - gamma 

        # Euler's method: new state = old state + rate * step size
        logS[t + 1] = logS[t] + dlogS * dt
        logI[t + 1] = logI[t] + dlogI * dt
        
    return S, I
\end{lstlisting}



%% \section{Observation noise and process noise}

%% A simulated (or real) epidemic time series has two distinct sources
%% of randomness, and it matters which one you mean.

%% \begin{itemize}
%%   \item \textbf{Process noise}: randomness in the \emph{true}
%%   dynamics themselves --- e.g.\ random fluctuations in the contact
%%   rate. This perturbs the actual trajectory $S_t,I_t,R_t$, step by
%%   step, as it is generated.
%%   \item \textbf{Observation noise}: randomness in how that true state
%%   is \emph{measured} --- underreporting, test error, sampling. This
%%   is added only once, after the true trajectory is simulated, to
%%   produce the ``data'' we pretend to observe.
%% \end{itemize}
%% From data alone, the two are often hard to tell apart --- a central
%% difficulty in fitting epidemic models, and part of the motivation for
%% the methods we turn to at the end of this chapter.

\section{A driven epidemic model}

Many infectious diseases are seasonal (e.g.~measles and influenza),
since transmission varies in time. We let $\beta$ vary
periodically in time:
\[
\beta_t = \beta_0\big(1+\beta_1\cos(2\pi t)\big),
\]
This is a \emph{driven}, or \emph{forced} system, with $t$ in years,
$\beta_0>0$ the baseline transmission rate, and $\beta_1\in[0,1)$ the
amplitude of seasonal forcing. Substituting into the demographic SIR
system:
\begin{align*}
\dot S &= \mu(1-S) - \beta_t\, S I, \\
\dot I &= \beta_t\, S I - (\gamma+\mu) I, \\
\dot R &= \gamma I - \mu R.
\end{align*}
%% This is a purely deterministic system, driven by a perfectly periodic
%% $\beta_t$, with no noise anywhere --- yet, for suitable parameters, it
%% produces recurring outbreaks that are not all the same size:

\begin{figure}[h]
\centering
\includegraphics[width=0.8\linewidth]{visualizations/driven_sir.pdf}
\caption{\label{fig:driven_sir}A biennial cycle: one large outbreak (red), one small (blue),
repeating exactly every two years ($\log$ scale --- the small
outbreaks are $\sim\!15\times$ smaller). Parameters: $\mu=0.02\,\text{yr}^{-1}$,
$\gamma=365/13\,\text{yr}^{-1}$ (13-day infectious period), $R_0=17$
(so $\beta_0=R_0(\gamma+\mu)\approx477.6\,\text{yr}^{-1}$),
$\beta_1=0.185$.}
\end{figure}

%% \section{Adding an incubation stage: the SEIR model}
%% The SIR model above has no incubation period: a susceptible becomes
%% infectious immediately on contact. For measles this is unrealistic ---
%% symptoms (and infectiousness) begin only $\sim\!8$--$12$ days after
%% exposure. We add an \emph{exposed} (infected but not yet infectious)
%% compartment $E$:
%% \begin{align*}
%% \dot S &= \mu(1-S) - \beta_t\, S I, \\
%% \dot E &= \beta_t\, S I - (\mu+a) E, \\
%% \dot I &= a E - (\mu+\gamma) I, \\
%% \dot R &= \gamma I - \mu R,
%% \end{align*}
%% where $1/a$ is the mean latent period.

\section{Units}
Units are crucial in simulations and modeling. We first decide on time
unit (e.g.~year, month, week) and then derive parameters
accordingly. See Table~\ref{tab:units} for examples.

\begin{table}[h]
\centering
\footnotesize
\begin{tabular}{l p{10.5cm}}
\hline
description & value \\
\hline
recovery rate &
$\gamma = \frac{1}{13}\,\text{day}^{-1} = \frac{1}{13}\left(\frac{\text{year}}{365}\right)^{-1}= \frac{365}{13}\,\text{year}^{-1}$ \\
%
%
%
birth rate & 
$\mu = \frac{1}{50}\,\text{year}^{-1} = \frac{1}{50}\big(365\,\text{day}\big)^{-1} = \frac{1}{18250}\,\text{day}^{-1}$ \\
%
%
%
contact rate &
$\kappa = 10\,\text{day}^{-1} =10\left(\frac{\text{year}}{365}\right)^{-1} = 3650\,\text{year}^{-1}$
\\ \hline
\end{tabular}
\caption{Converting parameter values between day and year units, by
treating day and year as algebraic symbols (day $=$ year$/365$).}
\label{tab:units}
\end{table}
